% GENERATED FILE. Edit canonical lessons, then run npm run generate:books.
% canonical-source-hash: fnv1a64:9272f771bbf5c68a
% canonical-lessons: KA-C159-suli, KA-C159-are, KA-C159-badisu, KA-C159-ruci-nodu, KA-C159-agi

\chapter{To peel, To grind, To serve, To taste, To chew}
\label{ch:ka-to-peel-to-grind-to-serve-to-taste-to-chew}
\hlchaptermodality{159}

\begin{chapteropening}
\textbf{By the end of this chapter:} \emph{I can say to peel, to grind, to serve, to taste and to chew.}

Complete the last lesson of chapter 159: I can say to peel, to grind, to serve, to taste and to chew.
\end{chapteropening}

\section[\texorpdfstring{\kn{ಸುಲಿ} (suli)}{ಸುಲಿ (suli)}]{\kn{ಸುಲಿ} (suli) — to peel}

\label{lesson:KA-C159-suli}



\begin{quote}
Before the new one: say the Kannada for to pay, then the Kannada for to bathe.
\end{quote}

\subsection*{You'll want to know: \kn{ಸುಲಿ}}
\textbf{\kn{ಸುಲಿ}} — \emph{suli} — \textquotedblleft{}to peel\textquotedblright{}.

\begin{grammarlens}[title={What you've built}]
One more word for this chapter.
\end{grammarlens}

\subsection*{Your turn}
\begin{itemize}
\raggedright
  \item \emph{Say it:} \emph{suli}
  \item \emph{Say it:} \emph{suli}, once more
  \item \emph{Say it:} say \emph{suli}
  \item \emph{Recall:} say the Kannada for to pay, then the Kannada for to bathe, then say \emph{suli} again
\end{itemize}

\begin{tcolorbox}[breakable,colback=teal!4,colframe=teal!35!black,title={Before you move on}]
What does \kn{ಸುಲಿ} mean? (\textquotedblleft{}To peel\textquotedblright{}.) Say it once more.
\end{tcolorbox}
\section[\texorpdfstring{\kn{ಅರೆ} (are)}{ಅರೆ (are)}]{\kn{ಅರೆ} (are) — to grind}

\label{lesson:KA-C159-are}



\begin{quote}
Before the new one: say the Kannada for to bathe, then the Kannada for to peel.
\end{quote}

\subsection*{You'll want to know: \kn{ಅರೆ}}
\textbf{\kn{ಅರೆ}} — \emph{are} — \textquotedblleft{}to grind\textquotedblright{}.

\begin{grammarlens}[title={What you've built}]
One more word for this chapter.
\end{grammarlens}

\subsection*{Your turn}
\begin{itemize}
\raggedright
  \item \emph{Say it:} \emph{are}
  \item \emph{Say it:} \emph{are}, once more
  \item \emph{Say it:} say \emph{are}
  \item \emph{Recall:} say the Kannada for to bathe, then the Kannada for to peel, then say \emph{are} again
\end{itemize}

\begin{tcolorbox}[breakable,colback=teal!4,colframe=teal!35!black,title={Before you move on}]
What does \kn{ಅರೆ} mean? (\textquotedblleft{}To grind\textquotedblright{}.) Say it once more.
\end{tcolorbox}
\section[\texorpdfstring{\kn{ಬಡಿಸು} (baḍisu)}{ಬಡಿಸು (baḍisu)}]{\kn{ಬಡಿಸು} (baḍisu) — to serve (food)}

\label{lesson:KA-C159-badisu}



\begin{quote}
Before the new one: say the Kannada for to peel, then the Kannada for to grind.
\end{quote}

\subsection*{You'll want to know: \kn{ಬಡಿಸು}}
\textbf{\kn{ಬಡಿಸು}} — \emph{baḍisu} — \textquotedblleft{}to serve (food)\textquotedblright{}.

\begin{grammarlens}[title={What you've built}]
One more word for this chapter.
\end{grammarlens}

\subsection*{Your turn}
\begin{itemize}
\raggedright
  \item \emph{Say it:} \emph{baḍisu}
  \item \emph{Say it:} \emph{baḍisu}, once more
  \item \emph{Say it:} say \emph{baḍisu}
  \item \emph{Recall:} say the Kannada for to peel, then the Kannada for to grind, then say \emph{baḍisu} again
\end{itemize}

\begin{tcolorbox}[breakable,colback=teal!4,colframe=teal!35!black,title={Before you move on}]
What does \kn{ಬಡಿಸು} mean? (\textquotedblleft{}To serve (food)\textquotedblright{}.) Say it once more.
\end{tcolorbox}
\section[\texorpdfstring{\kn{ರುಚಿ} \kn{ನೋಡು} (ruci nōḍu)}{ರುಚಿ ನೋಡು (ruci nōḍu)}]{\kn{ರುಚಿ} \kn{ನೋಡು} (ruci nōḍu) — to taste}

\label{lesson:KA-C159-ruci-nodu}



\begin{quote}
Before the new one: say the Kannada for to grind, then the Kannada for to serve.
\end{quote}

\subsection*{You'll want to know: \kn{ರುಚಿ} \kn{ನೋಡು}}
\textbf{\kn{ರುಚಿ} \kn{ನೋಡು}} — \emph{ruci nōḍu} — \textquotedblleft{}to taste\textquotedblright{}.

\begin{grammarlens}[title={What you've built}]
One more word for this chapter.
\end{grammarlens}

\subsection*{Your turn}
\begin{itemize}
\raggedright
  \item \emph{Say it:} \emph{ruci nōḍu}
  \item \emph{Say it:} \emph{ruci nōḍu}, once more
  \item \emph{Say it:} say \emph{ruci nōḍu}
  \item \emph{Recall:} say the Kannada for to grind, then the Kannada for to serve, then say \emph{ruci nōḍu} again
\end{itemize}

\begin{tcolorbox}[breakable,colback=teal!4,colframe=teal!35!black,title={Before you move on}]
What does \kn{ರುಚಿ} \kn{ನೋಡು} mean? (\textquotedblleft{}To taste\textquotedblright{}.) Say it once more.
\end{tcolorbox}
\section[\texorpdfstring{\kn{ಅಗಿ} (agi)}{ಅಗಿ (agi)}]{\kn{ಅಗಿ} (agi) — to chew}

\label{lesson:KA-C159-agi}



\begin{quote}
Before the new one: say the Kannada for to serve, then the Kannada for to taste.
\end{quote}

\subsection*{You'll want to know: \kn{ಅಗಿ}}
\textbf{\kn{ಅಗಿ}} — \emph{agi} — \textquotedblleft{}to chew\textquotedblright{}.

\begin{grammarlens}[title={What you've built}]
One more word for this chapter.
\end{grammarlens}

\subsection*{Your turn}
\begin{itemize}
\raggedright
  \item \emph{Say it:} \emph{agi}
  \item \emph{Say it:} \emph{agi}, once more
  \item \emph{Say it:} say \emph{agi}
  \item \emph{Recall:} say the Kannada for to serve, then the Kannada for to taste, then say \emph{agi} again
\end{itemize}

\begin{tcolorbox}[breakable,colback=teal!4,colframe=teal!35!black,title={Before you move on}]
What does \kn{ಅಗಿ} mean? (\textquotedblleft{}To chew\textquotedblright{}.) Say it once more.
\end{tcolorbox}
